\chapter{Conclusão}

O desenvolvimento do Sistema Web Administrativo para a Confeitaria Bom Gosto representa uma proposta para aprimorar a organização e a eficiência das atividades internas da empresa. A solução foi planejada para centralizar as informações relacionadas aos pedidos, produtos, estoque, ingredientes, produção e controle financeiro, contribuindo para reduzir a dependência de registros manuais e facilitar o acompanhamento das operações.

Nesta etapa do projeto, a estruturação do ambiente de desenvolvimento, a implementação da prova de conceito (POC) do módulo de autenticação e a configuração do banco de dados relacional representam avanços importantes para a construção da aplicação. Esses resultados estabelecem uma base para o desenvolvimento das regras de negócio e a integração das funcionalidades previstas.

Embora o sistema ainda necessite da implementação e validação de seus módulos, a análise realizada permitiu organizar os requisitos e definir as principais funcionalidades necessárias para atender às demandas administrativas da confeitaria. Dessa forma, o projeto apresenta uma base para a continuidade do desenvolvimento e para a futura utilização da aplicação no gerenciamento do negócio.


\section{Considerações Finais}

A tecnologia pode contribuir significativamente para a organização e o desenvolvimento de pequenos negócios, especialmente por meio da automatização de processos e da centralização das informações. No contexto da Confeitaria Bom Gosto, o sistema proposto busca facilitar o gerenciamento das atividades administrativas, proporcionando maior controle sobre os pedidos, os produtos, os ingredientes e os custos de produção.

A estruturação do ambiente de desenvolvimento utilizando a abordagem de \textit{monorepo}, juntamente com a implementação da prova de conceito do módulo de autenticação e a configuração do banco de dados relacional, demonstra o avanço na preparação da infraestrutura da aplicação. Entretanto, ainda é necessário implementar e integrar os demais módulos para que as funcionalidades previstas possam ser utilizadas de forma completa.

Após a implementação e a validação dos componentes, espera-se que o sistema contribua para tornar os processos mais organizados, facilitar o acesso às informações e auxiliar a proprietária na tomada de decisões relacionadas ao funcionamento da confeitaria.


\section{Dificuldades Encontradas}

Durante o desenvolvimento do projeto até a presente etapa, as principais dificuldades estiveram relacionadas à estruturação do ambiente de desenvolvimento, à configuração da infraestrutura e à definição de práticas de segurança e testes automatizados. A utilização de ferramentas como Docker, Angular e Django, além dos recursos de padronização de código, exigiu aprendizado e adaptação por parte da equipe para estabelecer um fluxo de trabalho adequado.

Outra dificuldade foi transformar os procedimentos administrativos da confeitaria em requisitos e regras de negócio que pudessem ser representados computacionalmente. A definição das funcionalidades relacionadas ao gerenciamento de pedidos, ao controle de estoque e inventário, à produção e às finanças exigiu a análise das necessidades da empresa e a organização dos processos que deverão ser contemplados pela aplicação.

Essas etapas foram importantes para estabelecer uma compreensão mais clara do funcionamento do negócio e orientar a implementação dos módulos previstos para o sistema.


\section{Trabalhos Futuros}

Como trabalhos futuros, prevê-se a continuidade do desenvolvimento dos módulos de Usuário, Pedido, Estoque, Produção e Inventário, com a realização de ajustes e validações conforme os requisitos definidos para o projeto. Também está prevista a implementação dos módulos de Agenda, Relatórios e Finanças, que ainda não foram desenvolvidos nesta etapa.
Após a implementação desses componentes, será necessário realizar a integração entre os módulos, validar os fluxos de funcionamento do sistema e executar testes para verificar se as funcionalidades atendem aos requisitos estabelecidos. Também poderão ser realizados ajustes na interface e nos mecanismos de segurança, visando melhorar a usabilidade, a confiabilidade e a proteção das informações armazenadas.
Por fim, a validação do sistema com base nas necessidades da proprietária permitirá identificar possíveis melhorias e orientar futuras evoluções da aplicação.